\cleardoublepage
\phantomsection
\pdfbookmark{Abstract}{Abstract}
\begingroup
\let\clearpage\relax
\let\cleardoublepage\relax
\let\cleardoublepage\relax

\chapter*{Abstract}


Over the past few decades, the use of mobile devices has increased drastically, partially replaces the usage of desktop-based software and web browsers. This shift has been mainly driven by the widespread adoption of smartphones, their increased computational capabilities, but also for the availability of integrated sensors such as GPS, gyroscopes and Bluetooth, whose, together allow richer and more context-aware user experiences. Consequently, by leveraging these technologies, Online Social Networks (OSNs) have been able to provide more advanced and sophisticated features, collecting larger amounts of more complex metadata, and thus increase platform interactions and daily engagement.\\


This thesis presents GeoMedia, an open-source framework designed for sharing multimedia content, such as text, images, audio and files, anchored to real-world geographic locations. By exploiting GPS sensors equipped on mobile devices, users can explore and publish digital content through an interactive map rendered by the application, creating a seamless integration between the digital and physical worlds. Moreover, the framework offers several features and covers various use cases, including restricting content visibility within a customizable distance range, enabling time-sensitive content or applying user-based restrictions for  content creation and visualization. GeoMedia also allows users to create sequential posts, forming an ordered chain for discovering posts and their related places. \\

% Furthermore, this thesis analyzes the motivations behind GeoMedia, its critical aspects and design considerations, describes its system architecture and implementation and compares different technological paradigms. Finally, it explores potential applications and outlines possible future developments of the framework.

% \vspace{1em}

% A previous version of this application was presented at the GoodIT 2025: International Conference on Information Technology for Social Good, held from September 3 to 5, 2025, in Antwerp, Belgium \cite{site:goodit}.

%\vfill

\endgroup

\vfill
